\chapter{PENELITIAN TERKAIT}

\section{Konsep Atensi Visual dan Prediksi Saliency}

% TODO: Jelaskan teori atau konsep pertama yang relevan dengan penelitian Anda.
% Gunakan referensi dari pustaka yang terpercaya.

Konsep \textit{saliency} dalam konteks atensi visual pertama kali diformalkan secara komputasional oleh \citet{itti1998model} melalui model \textit{bottom-up} yang terinspirasi dari arsitektur neuronal sistem visual primata. Model ini menggabungkan fitur intensitas, orientasi, dan warna dari sebuah citra ke dalam satu peta topografis (\textit{saliency map}) yang merepresentasikan seberapa besar kemungkinan suatu lokasi menarik atensi visual manusia secara alami tanpa dipengaruhi oleh tugas atau tujuan kognitif tertentu \citep{itti1998model}. Model-model generasi awal ini menjadi fondasi konseptual bagi seluruh perkembangan riset prediksi atensi visual selanjutnya, termasuk model berbasis deep learning yang lebih modern dan dibahas pada subbab 2.2.

Istilah \textit{``Saliency''} pada riset ini merujuk secara spesifik pada prediksi atensi visual/gaze, yaitu upaya memodelkan ke mana mata manusia secara alami akan tertarik ketika melihat sebuah antarmuka atau citra. Istilah ini perlu dibedakan dari penggunaan kata \textit{``saliency''} pada bidang riset yang berbeda, yaitu \textit{explainability} atau \textit{Explainable AI} (XAI), dimana teknik seperti \textit{Gradient-weight Class Activation Mapping} (Grad-CAM) digunakan untuk menghasilkan peta yang menyoroti wilayah citra mana yang paling memengaruhi keputusan klasifikasi sebuah model \textit{deep learning}, misalnya menjelaskan bagian citra X-ray mana yang mendorong model mendiagnosis suatu penyakit \citep{selvaraju2017grad}. Kedua bidang ini sama-sama menghasilkan keluaran visual berupa peta panas (heatmap), namun secara fundamental berbeda berdasarkan apa yang sedang dimodelkan. Peta atensi visual memodelkan persepsi dan perhatian mata manusia, sedangkan peta \textit{explainability} memodelkan dasar pengambilan keputusan sebuah model AI. Kerancuan antara dua istilah ini merupakan kekeliruan umum dalam penelusuran literatur di bidang ini, sehingga penelitian ini secara konsisten menggunakan istilah \textit{visual attention prediction} atau \textit{gaze/eye-tracking saliency} untuk menghindari ambiguitas tersebut.

Seiring perkembangan \textit{deep learning}, pendekatan prediksi atensi visual bergeser dari model \textit{bottom-up} berbasis fitur tangan (\textit{hand-crafted features}) seperti \citet{itti1998model} menuju model yang dilatih pada data eye-tracking berskala besar untuk mempelajari pola atensi manusia secara \textit{data-driven}. Model-model inilah yang menjadi fokus pembahasan pada subbab berikutnya.

Atensi visual manusia umumnya dipahami dipengaruhi oleh dua kelompok faktor: (1) faktor \textit{bottom-up} yang bersumber dari karakteristik visual stimulus (mis. kontras, warna, dan orientasi, sebagaimana dimodelkan oleh \citealp{itti1998model}) dan (2) faktor \textit{top-down} yang bersumber dari tujuan, tugas, atau pengetahuan pengamat \citep{itti2001computational}. Model prediksi saliency yang dibahas pada subbab berikutnya dilatih dan dievaluasi pada data atensi yang pada umumnya dikumpulkan tanpa tugas spesifik (\textit{free-viewing}; mis. UEyes, \citealp{jiang2023ueyes}; SALICON, \citealp{jiang2015salicon}), karena pengamatan bebas paling lazim dipakai dalam pemodelan saliency \citep{bylinskii2019what}, sehingga keluarannya terutama merepresentasikan komponen yang dapat dipelajari dari tampilan visual semata. Konsekuensinya, pada antarmuka yang digunakan untuk menyelesaikan tugas tertentu seperti openIMIS, peta atensi prediktif tidak dapat disamakan dengan pola atensi pengguna yang sedang bertugas. Studi \textit{eye-tracking} pada laman web melaporkan bahwa karakteristik atensi visual bergantung pada tugas yang dikerjakan pengguna \citet{buscher2009what}, sehingga perbedaan antara pengamatan bebas dan pengamatan berbasis tugas perlu diperhitungkan dalam menafsirkan hasil.


\section{Model Prediksi Saliency Berbasis Deep Learning}

Sejak tahun 2014, \textit{transfer learning} dari jaringan yang telah dilatih pada \textit{classification task} citra berskala besar (mis. ImageNet) menjadi pendorong utama peningkatan akurasi model prediksi saliency, menggantikan pendekatan \textit{bottom-up} berbasis fitur tangan yang dibahas pada subbab sebelumnya. Empat model berikut dipilih sebagai representasi \textit{state of the art} yang relevan bagi penelitian ini, karena seluruhnya tersedia sebagai model \textit{pretrained} yang dapat digunakan langsung tanpa pelatihan ulang. 
 
\begin{enumerate}
    \item \textbf{DeepGaze IIE} \citep{linardos2021deepgaze} dibangun di atas arsitektur DeepGaze II dengan mengganti \textit{backbone} VGG-19 menjadi kombinasi ensemble dari beberapa \textit{backbone} ImageNet yang lebih modern (mis. ResNet-50), digabungkan melalui jaringan \textit{readout} non-linear. Model ini dilatih awal (\textit{pretraining}) pada dataset SALICON (10.000 citra) kemudian di-\textit{fine-tune} pada dataset MIT1003 (citra natural). Karakteristik yang membedakan DeepGaze IIE dari pendahulunya adalah fokus eksplisit pada kalibrasi prediksi, baik pada domain data latih maupun pada domain yang belum pernah dilihat sebelumnya (\textit{out-of-domain}), sebuah properti yang secara langsung relevan dengan konteks penelitian ini, dimana model akan diterapkan pada domain openIMIS yang sama sekali tidak pernah menjadi bagian dari data latihnya.

    \item \textbf{TranSalNet} \citep{lou2022transalnet} mengintegrasikan \textit{encoder} CNN (tersedia dalam dua varian: ResNet-50 pada \textit{TranSalNet\_Res}, atau DenseNet-161 pada \textit{TranSalNet\_Dense}) dengan encoder Transformer yang menerapkan mekanisme \textit{multi-head-self-attention}, digabungkan melalui decoder CNN untuk merekonstruksi peta saliency akhir. Penggabungan ini dirancang agar model dapat menangkap fitur visual lokal (lewat CNN) sekaligus dependensi spasial jarak jauh antar elemen pada citra (lewat Transformer), karakteristik ini berpotensi relevan untuk antarmuka administratif seperti openIMIS yang padat elemen dan tabel data. Model ini dilatih pada dataset SALICON. TranSalNet\_Res dipilih sebagai model utama penelitian karena risiko kebutuhan VRAM yang lebih rendah dibanding TranSalNet\_Dense, dengan performa yang tetap kompetitif menurut evaluasi pada paper aslinya.

    \item \textbf{UMSI (Unified Model of Saliency and Importance)} \citep{fosco2020predicting} secara konseptual sedikit berbeda dari dua model sebelumnya yang memprediksi saliency berbasis fiksasi mata, UMSI memprediksi \textit{visual importance}, dilatih pada dataset Imp1k yang memuat 1.000 desain grafis (mencakup laman web, poster film, antarmuka perangkat mobile, infografis, dan iklan) dengan keterangan kepentingan visual yang dikumpulkan secara \textit{crowdsourced} melalui antarmuka anotasi manual, bukan data \textit{eye-tracking} fisik. UMSI juga dilengkapi modul klasifikasi otomatis untuk mengenali jenis desain input sebelum melakukan prediksi. Model ini penting disebut karena temuan dari \citet{jiang2023ueyes} yang menemukan bahwa ketika UMSI dilatih ulang secara khusus menggunakan data \textit{eye-tracking} UI dari UEyes, akurasinya (AUC) meningkat signifikan dari 0,778 menjadi 0,878, hal ini menjadi sebuah bukti kuantitatif bahwa pelatihan spesifik-domain memberikan pengaruh besar terhadap akurasi prediksi saliency pada antarmuka digital, dan menjadi dasar argumen keterbatasan pada pendekatan \textit{zero-shot} yang diusulkan penelitian ini.
    
    \item \textbf{SUM (Saliency Unification through Mamba)} \citep{hosseini2025sum} adalah model terbaru yang diusulkan, mengintegrasikan arsitektur \textit{Mamba} (\textit{State Space Model}) dengan struktur U-Net serta modul \textit{Conditional Visual State Space} (C-VSS) untuk beradaptasi terhadap berbagai jenis citra dalam satu model tunggal. SUM dilatih secara gabungan pada enam dataset benchmark berskala besar yang mencakup citra natural, laman web, dan citra \textit{e-commerce}. Berdasarkan Tabel 1 pada paper aslinya, salah satu dari enam dataset pelatihan SUM adalah UEyes (dicantumkan sebagai \textit{``U-EYE''}, kategori \textit{Web page}, 1.979 citra dengan 1.583 di antaranya digunakan untuk training), yaitu dataset yang sama yang digunakan sebagai benchmark validasi umum pada pendekatan \textit{zero-shot} tanpa \textit{fine-tuning} yang diusulkan pada Bab 1. Dengan demikian, pengujian SUM pada UEyes tidak dapat dianggap sebagai generalisasi zero-shot yang murni, karena model tersebut telah dilatih menggunakan sebagian data dari dataset yang sama. Temuan ini menjadi pertimbangan tambahan mengapa TranSalNet\_Res, yang tidak dilatih pada UEyes, dipandang sebagai pilihan yang tepat untuk mengukur kredibilitas umum model secara independen pada penelitian ini.
\end{enumerate} 

Secara umum, keempat model yang disebutkan di atas menunjukan pola konsisten di mana seluruhnya mamanfaatkan \textit{transfer learning} dari representasi visual berskala besar (baik dari ImageNet maupun data saliency gabungan), dan tidak dievaluasi atau dilatih khusus pada antarmuka sistem informasi kesehatan administratif seperti openIMIS, sehingga mengonfirmasi kembali kesenjangan yang diidentifikasi pada Bab 1.


\section{Benchmark Saliency pada Antarmuka Digital}

Keempat model yang dibahas pada subbab sebelumnya dievaluasi dan dibandingkan menggunakan dataset benchmark publik. Salah satu dataset yang relevan pagi penelitian ini adalah UEyes \citep{jiang2023ueyes}, sebuah dataset \textit{eye-tracking} berskala besar yang secara khusus dikumpulkan untuk antarmuka digital, bukan citra natural seperti dataset-dataset pendahulunya (SALICON, MIT1003, CAT2000). UEyes terdiri atas 1.980 tangkapan layar antarmuka, dengan masing-masing 495 citra untuk empat kategori UI, yaitu laman web (\textit{webpage}), antarmuka desktop, antarmuka \textit{mobile}, dan poster. Data atensi dikumpulkan dari 62 partisipan menggunakan perangkat \textit{eye-tracker in-lab} beresolusi tinggi (Gazepoint GP3, 60Hz), bukan metode proksi tidak langsung seperti pelacakan kurso mouse atau estimase gaze berbasis webcam yang digunakan pada beberapa dataset saliency sebelumnya. Dataset ini menyediakan peta saliency pada tiga durasi pengamatan (1,3, dan 7 detik) beserta data \textit{scanpath} dan metadata jenis citra untuk tiap sampel.

Evaluasi akurasi model prediksi saliency umumnya dilakukan menggunakan beberapa metrik standar, yang masing-masing mengukur aspek kesesuaian berbeda antara peta saliency hasil prediksi dan data \textit{ground truth}, serta memiliki karakteristik, kelebihan, dan kelemahan yang berbeda satu sama lain \citep{bylinskii2019what}

\begin{enumerate}
    \item \textbf{Correlation Coefficient (CC)} mengukur korelasi linear (Pearson) antara peta prediksi dan peta kepadatan atensi acuan, dengan memperlakukan kedua peta sebagai variabel acak yang dibandingkan secara piksel-demi-piksel. Nilai CC berkisar antara -1 hingga 1, dengan nilai mendekati 1 menunjukan kesesuaian yang semakin baik. Karakteristik penting dari CC adalah sifatnya yang simetris, yaitu memberi bobot penalti yang setara terhadap kesalahan \textit{false positive} (memprediksi atensi tinggi pada area yang sebenarnya tidak diperhatikan) maupun \textit{false negative} (gagal memprediksi atensi tinggi pada area yang sebenarnya diperhatikan)

    \item \textbf{Normalized Scanpath Saliency (NSS)} mengukur nilai rata-rata peta prediksi (setelah dinormalisasi menjadi rata-rata nol dan standar deviasi satu) secara spesifik pada titik-titik lokasi fiksasi aktual pengguna, bukan pada keseluruhan peta. Karena hanya dihitung pada lokasi fiksasi diksrit, NSS lebih sensitif terhadap presisi lokasi dibanding CC. Nilai NSS positif menunjukan kesesuaian prediksi yang lebih baik dari tebakan acak, dan nilai negatif menunjukan prediksi yang justru berlawanan dengan lokasi atensi sebenarnya.

    \item \textbf{Kullback-Leibler Divergence (KLD)} mengukur seberapa besar perbedaan antara dua distribusi probabilitas, yaitu peta prediksi dan peta acuan yang telah dinormalisasi agar jumlah nilainya sama dengan satu. Berbeda dengan CC dan NSS, nilai KLD yang semakin kecil (mendekati nol) menunjukan kesesuaian yang semakin baik. KLD bersifat asimetris sehingga arah perhitungannya menentukan jenis kesalahan yang lebih dihukum. Bila divergensi dihitung dari peta acuan terhadap peta prediksi, sebagaimana konvensi yang dibahas \citet{bylinskii2019what}, penalti lebih besar dijatuhkan pada negatif palsu (area yang diperhatikan pengguna tetapi tidak diprediksi). Implementasi evaluasi resmi UEyes yang diadaptasi pada penelitian ini menghitung divergensi dari peta prediksi terhadap peta acuan, sehingga penalti lebih besar dijatuhkan pada positif palsu (atensi diprediksi pada area yang tidak diperhatikan). Arah yang dipakai dinyatakan eksplisit pada Persamaan (3.7) di Bab 3.

    \item \textbf{Area Under the ROC Curve (AUC)} memperlakukan evaluasi sebagai persoalan klasifikasi biner, yaitu membedakan antara titik fiksasi dan bukan-fiksasi berdasarkan ambang nilai (\textit{treshold}) pada peta prediksi. Nilai AUC berkisat antara 0,5 (setara tebakan acak) hingga 1 (klasifikasi sempurna). Karena berbasis pada peringkat (\textit{ranking}) nilai piksel dan bukan nilai absolutnya, AUC relatif tidak sensitif terhadap skala atau normalisasi peta prediksi, namun beberapa variannya (mis. \textit{shuffled AUC}) perlu digunakan secara hati-hati untuk mengoreksi bias sentral (\textit{center bias}) pada data fiksasi manusia.
\end{enumerate}


Karena masing-masing metrik menangkap aspek kesesuaian yang berbeda, baik dari sisi simetrik/asimetris penalti kesalahan maupun representasi \textit{ground truth} yang digunakan (peta kepadatan kontinu untuk CC dan KLD, titik fiksasi diskrit untuk NSS dan AUC), penggunaan kombinasi beberapa metrik secara bersamaan lebih komprehensif dibandingkan mengandalkan satu metrik tunggal untuk menyimpulkan akurasi sebuah model \citep{bylinskii2019what}.

Sebagaimana telah disinggung pada subbab 2.2, temuan dari \citet{jiang2023ueyes} menunjukan bahwa pelatihan ulang model UMSI secara khusus menggunakan data UEyes meningkatkan akurasinya secara signifikan dibandingkan performa model tersebut saat belum dilatih pada domain antarmuka digital. Temuan ini mengindikasikan bahwa antarmuka digital memiliki karakteristik visual yang cukup berbeda dari citra natrual, sehingga akurasi model saliency generik tidak dapat diasumsikan berlaku sama baiknya pada domain antarmuka.

Karena UEyes merupakan dataset eye-tracking berskala besar dengan kategori yang mencakup antarmuka desktop dan laman web, yaitu dua kategori yang secara struktural paling mendekati antarmuka openIMIS (dibandingkan kategori mobile atau poster), dataset ini dipilih sebagai benchmark validasi kredibilitas umum pada penelitian ini, sebagaimana telah dijelaskan pada Bab 1.  

\section{Vision-Language Model Zero-shot untuk Prediksi Atensi Visual UI}

Selain model saliency konvensional yang dibahas pada subbab 2.2, pendekatan alternatif yang mulai dieksplorasi pada 2026 adalah pemanfaatan \textit{Vision-Language Model} (VLM) seara \textit{zero-shot} untuk memprediksi atensi visual pada antarmuka digital, tanpa pelatihan khusus pada data saliency sama sekali. Penelitian paling relevan pada arah ini adalah UIGaze \citep{song2026uigaze}, yang secara khusus meneliti seberapa dekat prediksi VLM dapat mendekati pola atensi visual manusia pada antarmuka digital, menggunakan dataset UEyes yang sama dengan yang digunakan pada penelitian ini sebagai \textit{ground truth}

Penelitian tersebut mengevaluasi sembilan VLM dari lima penyedia (OpenAI, Google, Anthropic, Qwen, dan ByteDance) melalui sebuah pipeline prediksi koordinat zero-shot. Setiap model diminta memprediksi 30 sampai 50 titik koordinat gaze dalam format JSON beserta nilai intensitasnya, yang kemudian diubah menjadi peta saliency kontinu melalui proses \textit{Gaussian blur}, lalu dibandingkan dengan data acuan UEyes menggunakan metrik CC, \textit{Similarity} (SIM), dan KLD. Tiga temuan utama yang dilaporkan pada penelitian ini:

\begin{enumerate}
    \item \textbf{Pertama}, seluruh VLM yang diuji menunjukan kesesuaian yang semakin baik terhadap data acuan berdurasi lebih panjant (7 detik) dibandingkan durasi singkat (1 detik), mengindikasikan bahwa VLM cenderung mendekati pola atensi eksploratif (setelah pengguna memiliki waktu untuk memahami keseluruhan antarmuka), bukan fiksasi awal yang bersifat reflektif. Model dengan kesesuaian tertinggi pada durasi 7 detik adalah GPT-5.4, dengan CC = 0,408, SIM = 0,503, dan KLD = 1,345.

    \item \textbf{Kedua}, tingkat kesesuaian antara prediksi VLM dan atensi manusia bervariasi tergantung karakteristik struktural antarmuka. Antarmuka desktop dengan hierarki visual yang jelas (navigasi, area konten, dan sidebar yang terpisah dengan baik) menghasilkan kesesuaian tertinggi (GPT-5.4 mencapai CC - 0,506 pada kategori ini), sedangkan antarmuka \textit{mobile} yang padat dengan elemen-elemen yang saling bersaing dalam ruang tampilan sempit menghasilkan kesesuaian yang jauh lebih rendah.

    \item \textbf{Ketiga}, kemampuan berinteraksi dengan antarmuka (mis. mengeklik, mengetik, menavigasi) dan kemampuan memprediksi atensi visual terbukti merupakan dua hal yang berbeda secara fundamental. UI-TARS 1.5, model yang dirancang khusus untuk operasi otomatis pada antarmuka, justru menunjukkan kesesuaian yang hampir nol dengan atensi manusia (CC = 0,086 pada durasi 7 detik, bahkan negatif pada kategori \textit{mobile}), menegaskan bahwa kemampuan "memahami ke mana pengguna akan melihat" tidak serta-merta dimiliki oleh model yang mahir "bertindak" pada antarmuka. 
\end{enumerate}

\citet{song2026uigaze} sendiri menyandingkan temuan \citet{jiang2023ueyes}, bahwa UMSI yang dilatih ulang khusus pada data UEyes mencapai AUC = 0,878, dengan hasil VLM terbaik pada penelitian mereka (GPT-5.4, CC - 0,408 dalam kondisi \textit{zero-shot} penuh tanpa pernah menyentuh data eye-tracking sama sekali). perbandingan ini melibatkan dua metrik yang berbeda (AUC dan CC), dan \citet{song2026uigaze} sendiri mengakui eksplisit bahwa keduanya tidak dapat dibandingkan secara langsung sebagai ukuran performa yang setara. Namun, perlu dicatat bahwa \citet{jiang2023ueyes} turut juga melaporkan skor CC untuk UMSI yang dilatih ulang pada UEyes, yaitu CC = 0,699. Dengan demikian perbandingan yang setara metrik tetap dapat disusun: CC = 0,699 untuk UMSI hasil pelatihan khusus domain, berbanding CC = 0,408 untuk VLM terbaik dalam kondisi \textit{zero-shot} penuh, selisih yang tetap mengindikasikan kesenjangan akurasi yang cukup besar akibat ketiadaan pelatihan domain-spesifik.

Keberadaan UIGaze memiliki implikasi penting bagi penelitian ini. Jika pendekatan VLM \textit{zero-shot} dipertimbangkan sebagai alternatif atau pembanding terhadap TranSalNet\_Res pada tahap penelitian selanjutnya, diferensiasi eksplisit terhadap UIGazw menjadi wajib, karena penelitian tersebut telah lebih dahulu menguji penekatan VLM \textit{zero-shot} pada dataset UEyes. Perbedaan utama yang membedakan penelitian ini dari UIGaze tetap terjaga, yaitu UIGaze mengevaluasi empat kategori UI umum pada UEyes (laman web, desktop, \textit{mobile}, poster), sedangkan penelitian ini berfokus secara spesifik pada antarmuka seistem informasi kesehatan administratif (openIMIS), yang masih menjadi topik terbatas dalam area penelitian saliency. 

\section{Usability dan Human-Computer Interaction pada Sistem Informasi Kesehatan}

Tantangan \textit{usability} pada sistem informasi kesehatan telah didokumentasikan pada berbagai studi sebelumnya, baik dari sisi hambatan implementasi maupun pendekatan yang diusulkan untuk mengatasinya.

\citet{verma2021openmrs} melakukan kajian retrospektif selama 15 tahun terhadap implementasi OpenMRS di berbagai negara berpendapatan rendah dan menengah. Studi ini mengiedentifikasi tiga hambatan implementasi yang paling umum dilaporkan, yaitu ketersediaan infrastruktur yang buruk (perangkat, listrik, dan konektivitas), kesulitan pengguna klinis bertransisi dari sistem berbasis kertas ke sistem digital akibat literasi teknologi informasi yang rendah, dan kesulitan menemukan staf teknis yang terampil di negara setempat. Berdasarkan temuan tersebut, \citet{verma2021openmrs} merekomendasikan perlunya sumber daya komunitas untuk mendorong \textit{usability} EHR, termasuk menciptakan versi OpenMRS yang siap pakai (\textit{shelf-ready}) dengan antarmuka pengguna yang terstandarisasi, serta menginetgrasikan \textit{Human-Centered Design} (HCD) ke dalam proses pengembangan sistem.

\citet{paton2022open}, mewakili \textit{IMIA Open Source Working Group}, mengusulkan tiga pilar utama untuk menilai dampak proyek \textit{open source} terhadap sistem kesehatan digital, yaitu \textit{resilience, accebility} dan \textit{equity}. Pilar \textit{accessibility} diilustrasikan secara khusus melalui pembahasan pembangunan ulang arsitektur antarmuka OpenMRS generasi baru (dikenal sebagai O3) menggunakan arsitektur \textit{Micro-frontend}, yang mengadopsi \textit{IBM Carbon Design System}, yaitu satu set pedomoan visual, pengalaman pengguna, dan kode yang tersedia secara terbuka, disertai pengujian pengguna secara ekstensif di berbagai negara. \citet{paton2022open} menyatakan bahwa pendekatan ini diharapkan ini diharapkan meningkatkan kemudahan pengguna serta mempermudah modifikasi dan kombinasi komponen antarmuka, sehingga membanutu mengatasi tantangan adaptasi sistem EHR serta pelatihan staf dengan pengalaman IT terbatas. Bagian yang secara khusus relevan bagi penelitian ini, \citet{paton2022open} juga menyebutkan openIMIS secara eksplisit sebagai salah satu proyek \textit{open source} (bersama OpenMRS, DHIS2, dan GNUHealth) yang telah menerapkan sistem desain dan kerangka kerja serupa yntuk mencapai tingkat usabilitas dan aksesibilitas yang tinggi bagi pasien maupun penyedia layanan kesehatan.

Sebagai contoh pendekatan evaluasi usability yang lebih konvensional, \citet{rick2015eyes} mengembangkan metodologi pengumpulan data eye-tracking secara in situ pada klinik rawat jalan untuk menganalisis bagaimana dokter berinteraksi dengan antarmuka EMR. Dengan menggunakan perangkat eye-tracker fisik yang dipasang pada 12 dokter di tiga klinik berbeda, penelitian tersebut berhasil mengumpulkan data dari 88 sesi kunjungan pasien, dengan total lebih dari 40 jam rekaman. Studi ini menunjukkan bahwa evaluasi \textit{usability} berbasis \textit{eye-tracking} konvensional mampu memberikan wawasan mendalam mengenai pola interaksi pengguna dengan antarmuka EMR, namun membutuhkan sumber daya logistik, perangkat keras khusus, dan waktu pengumpulan data yang besar untuk mengevaluasi satu sistem antarmuka saja.

Secara keseluruhan, studi-studi pada subbab ini menunjukkan bahwa tantangan \textit{usability} pada sistem informasi kesehatan telah diakui secara luas, dan berbagai pendekatan telah diusulkan untuk mengatasinya, mulai dari integrasi HCD ke proses pengembangan, adopsi \textit{design system} terstandardisasi, hingga evaluasi langsung melalui \textit{eye-tracking} konvensional. Meskipun demikian, ketiga pendekatan tersebut berfokus pada proses pengembangan dan evaluasi usabilitas secara umum, dan belum ada yang secara spesifik memanfaatkan model prediksi atensi visual berbasis \textit{deep learning} sebagai pendekatan analisis awal untuk mengidentifikasi risiko desain antarmuka pada sistem informasi kesehatan, khususnya pada antarmuka administratif seperti openIMIS.

\section{Pengodean Berbasis Dokumen dan Reliabilitas Pengode}

Pelabelan elemen informasi utama pada penelitian ini merupakan pengodean (\textit{coding}) terhadap sumber dokumenter, yaitu dokumentasi, materi pelatihan, dan video demonstrasi, menggunakan buku kode (\textit{codebook}) yang mendefinisikan kriteria secara tertulis. Buku kode yang terdokumentasi bertujuan membuat keputusan pengode dapat ditelusuri dan diulang oleh pihak lain. \citet{macqueen1998codebook} menjelaskan bahwa buku kode terstruktur menyediakan kerangka yang stabil bagi analisis data tekstual dan bahwa fitur-fitur tertentu pada buku kode dapat meningkatkan kesepakatan antar-pengode.

Konsistensi pengodean umumnya dinilai dengan mengukur kesepakatan antar-pengode (\textit{inter-rater}) atau antar-sesi oleh pengode yang sama (\textit{intra-rater}), salah satunya menggunakan koefisiane kappa yang mengoreksi kesepakatan terhadap kebetulan \citep{cohen1960coefficient}. \citet{hallgren2012computing} menguraikan pemilihan statistik, perhitungan, dan pelaporan kappa untuk data nominal. Rumus dan prosedur yang digunakan dijelaskan pada subbab 3.3.2. Perlu dicatat bahwa kesepakatan yang tinggi menunjukkan keterulangan keputusan, bukan kebenarannya. Pada penelitian ini tidak ada validasi oleh pakar domain, sehingga hasil pelabelan dinyatakan eksploratif.



\section{Kesenjangan Penelitian}

Berdasarkan uraian pada subbab 2.1 hingga 2.5, dapat disusun isntesis mengenai posisi penelitian ini terhadap penelitian-penelitian terkait yang telah di bahas.

Model prediksi saliency berbasis \textit{deep learning} (subbab 2.2) merupakan bidang riset yang telah matang, dengan sejumlah model pra-latih yang tersedia secara terbuka (DeepGaze IIE, TranSalNet, UMSI, SUM), namun seluruhnya dikembangkan dan dievaluasi pada dataset citra natural atau antarmuka digital secara umum, tanpa mempertimbangkan karakteristik khusus antarmuka sistem indormasi kesehatan. Dataset UEyes (subbab 2.3) memang secara khusus dikumpulkan untuk antarmuka digital, namun kategorinya (\textit{webpage}, desktop, \textit{mobile}, poster) bersifat umum dan tidak mencakup kategori antarmuka sistem kesehatan administratif atau klinis secara spesifik. Pendekatan VLM \textit{zero-shot} (subbab 2.4) telah diuji secara sistematis pada kategori UI umum melalui UIGaze, namun juga tidak menyentuh domain kesehatan. Sementara itu, penelitian pada ranah \textit{usability} sistem informasi kesehatan (subbab 2.5) secara konsisten mengangkat pentingnya evaluasi antarmuka yang baik, bahkan menyebutkan openIMIS secara eksplisit sebagai salah satu sistem yang relevan, namun pendekatan yang digunakan terbatas pada integrasi \textit{Human-Centered Design}, adopsi sistem desain terstandarisasi, atau evaluasi \textit{eye-tracking} konvensional yang memerlukan partisipan dan perangkat keras khusus, tanpa memanfaatkan model prediksi atensi visual berbasis \textit{deep learning} sebagai pendekatan analisis.

Dengan demikian, kesenjangan penelitian yang teridentifikasi adalah belum ditemukannya studi yang menerapkan model prediksi atensi visual berbasis \textit{deep learning}, khususnya dalam kondisi \textit{zero-shot} tanpa pelatihan ulang menggunakan data domain-spesifik, pada antarmuka sistem informasi kesehatan administratif, khususnya openIMIS. Kesenjangan ini konsisten ditemukan pada seluruh kategori literatur yang dibahas pada subbab 2.1 hingga 2.5, baik dari model (subbab 2.2), dataset benchmark (subbab 2.3), pendekatan VLM (subbab 2.4), maupun studi \textit{usability} sistem kesehatan (subbab 2.5).

Penelitian ini diposisikan untuk mengisi kesenjangan tersebut dengan menerapkan model prediksi atensi visual pra-latih secara \textit{zero-shot} pada antarmuka openIMIS, menggunakan UEyes sebagai benchmark validasi kredibilitas umum, dan membandingkan peta atensi yang dihasilkan dengan elemen informasi utama yang dilabeli secara sistematis dari sumber resmi openIMIS, sebagaimana telah dijelaskan pada Bab 1.

% TODO: Sesuaikan dengan aspek spesifik penelitian Anda.

